\newpage
\section{Modular Arithmetic}

We define an equivalence relation 

\[
    a \cong b \mod n \iff n | (a - b) \land  a,b \in \Z
\]

Means that both have the same residual when divided by \(n\). 

\subsection{Residual-class  and Residual-class Rings}

Because we have an equivalence relation we can classify the partitions of our set \(\Z\) into classes which we will 
call \emph{Residual-classes}. 

\textbf{Example:}

\[
    [13]_7 + \{m \in \Z \mid m \cong 13 \mod 7\}
\]

The set of all Residual-classes can be written as \(\Z / n \Z\) where \(n\) is the modulo. 


\subsection{Properties}

\emph{Addition}

\[
    (a + b) \mod n = (a \mod n + b \mod n) \mod n
\]

\emph{Multiplication}

\[
    (a * b) \mod n = (a \mod n * b \mod n) \mod n
\]

\emph{Exponentiation}

\[
    (a^m) \mod n = ({(a \mod n)}^n) \mod n
\]

\emph{Division}

\[
    \left(\frac{a}{b}\right) \mod n \ne  \frac{a \mod n}{b \mod n}
\]

\subsection{General Division Rules}

\subsubsection{Sum Rules}

If the number \(a\) can be expressed as a sum and \(b\) is a common divisor of all the terms, 
then \(b\) is also a divisor of \(a\).

\textbf{Example:} 

\( 742 = 700 + 42 \). Then 7, as a common divisor of 700 and 42, is a divisor of 742.

\subsubsection{Product Rule}

If the number \(a\) can be written as a product of two factors and \(b\) is a divisor of one of these 
factors, then \(b\) is also a divisor of \(a\).

\textbf{Example:}

\( 1500 = 15 \cdot 100 \). Then 3 and 5 are divisors of 15, and 2, 4, 5, 10, 25, and 50 are divisors of 
100—thus, all are divisors of 1500.

\subsubsection{Divisor Pair Rules}

If the number \(a\) has \(b\) as a divisor, then \( \frac{a}{b} \) is also a divisor of \(a\).

\textbf{Example:}

3 is a divisor of 315. Then \( 105 = \frac{315}{3} \) is also a divisor of 315.

\subsubsection{Divisor Rule}

If the number \(a\) has \(b\) as a divisor, then every divisor of \(b\) is also a divisor of 
\(a\).

\textbf{Example:}

45 is a divisor of 450. Then 3, 5, 9, and 15 (divisors of 45) are also divisors of 450.

\subsubsection{Difference Rule}

If the number \(a\) can be expressed as a difference and \(b\) is a common divisor of the minuend 
and subtrahend, then \(b\) is also a divisor of \(a\). \\

\textbf{Example:}

\( 686 = 700 - 14 \). Then 7, as a common divisor of 700 and 14, is a divisor of 686.

\subsubsection{Relative Prime Rule}

If the number \(a\) has two relatively prime numbers \(b\) and \( c \) as divisors, then their 
product \( b \cdot c \) is also a divisor of \(a\). \\

\textbf{Example:}

1540 has the relatively prime numbers 7 and 11 as divisors. Then \( 77 = 7 \cdot 11 \) is also a divisor 
of 1540.

\subsection{Using Modular Arithmetic to find the n digit of a big number}

What is the last digit of the number \(13^{2023}\)

Note that \(13 \mod 10 = 3\) and  \(13^2 \mod 10 =  ((13 \mod 10)  (13\mod 10)) \mod 10 = 9 \)

We can continue this process, and we will notice a pattern. After some iteration we will now that

\[(13^{2020} 13) \mod 10 = 7\]

\subsection{Gaussian Brackets}

The Gaussian brackets refer to two closely related functions used to round real numbers to integers 
in specific ways:

\emph{Floor Function}

The floor function, denoted by \(\lfloor x \rfloor\), gives the greatest integer less than or equal 
to a real number \(x\).

\textbf{Example:}

\(\lfloor 3.7 \rfloor = 3\), \(\lfloor -1.2 \rfloor = -2\)

\emph{Ceiling Function} 
    
The ceiling function, denoted by \(\lceil x \rceil\), gives the smallest integer greater than or 
equal to a real number \(x\).

\textbf{Example:}

\(\lceil 3.7 \rceil = 4\), \(\lceil -1.2 \rceil = -1\)

These functions are particularly useful: in number theory, computer science 
and rounding operations, where precise integer bounds of real numbers are required.

\subsection{Residual Class Rind modulo \(n\)}

Given \(n \in \N\). We define the relation \(\sim_n\) for \(\Z\).

\[
    a \sim b \iff n mid (a - b)
\]

Also, \(a / n\) and \(b / n\) have the same residual 

\subsection{Equivalence Class of \(\Z\)} 

For \(i, n \in \Z\) 

\[
    \overline{i} := i + n \Z = \{ i + na \mid a \in \Z\}
\]

and 

\[
    \Z_N = \Z / n \Z = \{\overline{0}:= n \Z, \overline{1}:= 1 + n \Z, \dots, \overline{n - 1} := n - 1 + n \Z \}    
\]

is the set of all equivalence relationships for \(\Z\) with respect to the relation 
\(k \sim l \iff n \mid (k - l)\). We define the addition and multiplication for \(\Z / n \Z\) using the 
representatives 

\[
    (a  + n \Z) + (b + n \Z) := (a + b) + n\Z
\]

\[
    (a  + n \Z) \cdot (b + n \Z) := ab + n\Z
\]

For all \(a, b \in \Z\). Therefore, 

\[
    \overline{a} + \overline{b} = \overline{a + b}, \quad \overline{a} \cdot \overline{b} = \overline{ab}
\]

